\scanpage{9}%
rer que pour $d$ assez grand, ce nombre est non nul (ce qui impliquera \deleted{\unsure{al}} alors\deleted{\unsure{,}} - et est sans doute équivalent - que le motif $E(Y)$ pour tous ces $Y$ est de \uline{niveau} maximum $n-1$, ou encore que le motif $E(Y)$ pour $Y$ générique est de niveau $n-1$). Sous la forme de (8.5.6)\deleted{\unsure{,}}, on voit que ceci est assuré dès que l'on sait que
\[
\dim E(Y,\underline O_Y)_{ss}\handwritten{>0}\deleted{\unsure{\longrightarrow+\infty}}\qquad\text{pour}\qquad d\handwritten{\longrightarrow}+\handwritten{\infty}\;,
\]
où le dim est pris bien entendu avec sa valeur générique. \deleted{\handwritten{On sait seulement (cf. n\textsuperscript{o} 7) que $\dim E(Y,\underline O_Y)_{ss}>0$ pour $d$ grand, et que $\varphi\ne0$ pour $d$ grand (donc $E(Y)$ de niveau $n-1$) dès que \unsure{l'on satisfait à|l'on sait que} $\dim(\mathrm{Tors}^{n+1}(X)\otimes k)_{ss}<p-1$.}}
% slanted
\margin{\handwritten{Or cela est établi au n\textsuperscript{o} 7}}

8.5.9. Le corps de base $k$ étant à nouveau quelconque (même pas \unsure{nécessairement} parfait), désignons, pour un degré $d$ donné, par $U_d$ l'ouvert de l'espace projectif des hypersurfaces de degré $d$ dans $\mathrm P^r_k$, formé des hypersurfaces dont l'intersection avec $X$ est \deleted{non singulière} \inserted{lisse} de dimension $n-1$, et telles que $E(Y,\underline O_Y)_{ss}$ soit de rang maximum. Alors lessuites exactes (8.5.1), pour les $Y$ provenant de points de $U_d$, proviennent d'une suite exacte de Modules localement libres \inserted{sur $U_d$} à opération $p$-linéaire semi-simple, ou de façon équivalente (cf 8.4.2) à une suite exacte de systèmes locaux \inserted{de vectoriels sur $\underline{\mathrm F}_p$} sur $U_d$ pour la top\deleted{\unsure{l}}ologie étale, \deleted{\unsure{formée d}} ou enfin, une fois choisi un point géométrique de $U_d$ (par exemple via une clôture \unsure{séparable|ségarable} du point générique) une suite exacte de représentations \inserted{de $G=\handwritten{\pi}_1(U_d)$} dans des vectoriels de dimension finie sur $\underline{\mathrm F}_p$. \deleted{\unsure{On}} On trouve un homomorphisme \inserted{\handwritten{$\varphi$}} du système local défini par les $E(Y)\otimes_W k'$ \deleted{dans} \inserted{\handwritten{\uline{sur}}} celui défini par les $E(Y,\underline O_Y)$, dont le\deleted{s} noyau \deleted{et le conoyau sont des} \inserted{\handwritten{est un}} système\deleted{s} loca\deleted{ux}\handwritten{l} \uline{constant} (i.e. provenant du corps de base $k$). Lorsque $k$ est fini, on peut supposer, quitte à passer à une extension finie de $k$, que les valeurs propres du frobenius relatif à $k$ opérant sur ce\deleted{s} dernier\deleted{s} \deleted{\unsure{(aux}} sont toutes égales à 1 (il suffit que les valeurs propres de frobenius relatif à $k$ opérant sur \deleted{les} $H^{n\deleted{\unsure{-1}}}(X,\underline O_X)_{ss}$ \deleted{et sur les deux membres de (8.5.3)} soient
% archive sheet number: 8
